% !TeX root = ../main.tex

\chapter{Baselines}\label{ch:baselines}

% Scope rule for this chapter: carry the baseline designs, the characterisation of the
% published comparator, and the reproducibility audit. Results ordered by research
% question belong in Chapter 6; interpretation belongs in Chapter 7.

This chapter states which floors and comparators the ladder in
\autoref{sec:protocol} is judged against, characterises the published comparator this
thesis reproduces, reports a reproducibility audit of that comparator alongside the
reproducibility properties of the implementation evaluated in this thesis, and states what
can and cannot be equalised in a head-to-head comparison between the two.

\section{The Baseline Ladder}\label{sec:baseline-ladder}

Six floors and comparators are read against every rung of the ladder described in
\autoref{sec:protocol}. Each is chosen to isolate a specific contrast rather than to
provide a single omnibus comparison.

The first floor uses age and sex alone, with no imaging information. Any arm that does not
exceed it is not using the graph at all. The second floor truncates each fold's own
validation subjects to their first visit and scores them at that fold's threshold; it
therefore never requires reading the held-out test set and asks whether a single scan
already carries the signal a trajectory model is meant to add. The third floor is a
logistic regression fitted on the principal components of the flattened connectivity-change
matrix, concatenated with visit count, total follow-up months, age and sex. It is a linear
floor: since it operates directly on raw connectivity change with no learned graph
representation, a longitudinal graph model that fails to exceed it is not adding
information beyond a linear read of the same change.

The fourth and fifth floors are both instances of the spatial-first \ac{GELSTM}
architecture introduced in \autoref{sec:spatial-first}, differing only in how the graph
encoder is initialised. One uses the pooled \ac{GAAE} checkpoint frozen, so the encoder
carries reconstruction pretraining over 3{,}700 unlabelled graphs before any classification
gradient reaches it. The other uses a randomly initialised encoder trained end to end
together with the recurrent head, so it learns only from the classification gradient on
248 labelled subjects. Comparing the pretrained-frozen \ac{GELSTM} against the
temporal-first arm that also uses a self-supervised encoder isolates the pooling-order
inversion under matched pretraining; comparing the randomly initialised \ac{GELSTM} against
the temporal-first arm with no pretraining on either side isolates the same inversion with
no pretraining confound. Reading either contrast alone would confound the architectural
inversion with the presence or absence of pretraining, which is why both floors are
carried through the ladder rather than one.

The sixth comparator is Brain-TokenGT, a published architecture evaluated here as a
reference point rather than as a controlled baseline; \autoref{sec:braintokengt}
characterises it and \autoref{sec:comparability} states why it cannot be equalised with the
other floors on every axis.

\section{Brain-TokenGT as a Published Comparator}\label{sec:braintokengt}

Brain-TokenGT \parencite{dong2023braintokengt} tokenises each visit's connectivity graph
into a sequence of node and edge tokens and applies cross-visit self-attention over the
resulting sequence, with a recurrent core that evolves the graph-convolution weights from
one visit to the next. This section characterises that recurrent core as released and under
the configuration used to obtain a trainable baseline in this thesis; the four other
respects in which the released implementation departs from the published architecture are
enumerated in \autoref{ch:app-implementation} together with the numeric settings of the
configuration adopted here.

The model definition evaluated in this thesis is the authors' released implementation,
carried over without changing the computation it performs. Two changes were necessary to
run it under the protocol used here: generalising a hardcoded region count and sequence
length to the values used on \ac{DELCODE}, and removing a hardcoded assignment of tensors
to a specific device. Neither changes what the model computes at the authors' original
configuration, which is confirmed by a regression test that copies the authors' released
weights into this implementation parameter for parameter and asserts that the two forward
passes agree at that configuration. The released training script trains on a single subject
with no held-out split, so training and evaluation for Brain-TokenGT use the same
cross-validation and held-out test harness applied to every other arm in this thesis. This
does not amount to reproducing the authors' reported result, which was obtained on a
different cohort, a different brain parcellation and the authors' own code path rather than
this one; no such reproduction is claimed here.

The recurrent core exhibits three distinct behaviours across the configurations evaluated
in this thesis. As released, the recurrent module's weights are never registered with the
optimiser, so they remain at their random initialisation throughout training. Under that
configuration the decision rule collapses to predicting conversion for every test subject,
which is an uninformative classifier rather than a trained one. A second configuration
registers the recurrent weights with the optimiser and trains them end to end, combined
with the other changes to the released implementation described in
\autoref{ch:app-implementation}; under it, training produces non-finite scores inside the
model's pooling layer and terminates within approximately ten epochs. A third configuration
retains those same changes and additionally applies a reduced learning rate and an added
weight-decay term to the recurrent parameters specifically; under it, training completes
without producing non-finite values, and this is the only one of the three configurations
for which a held-out evaluation is reported anywhere in this thesis. That a run completes
without diverging is a statement about training stability and is addressed by this
configuration; \autoref{sec:reproducibility-audit} addresses a separate question, namely
whether repeated runs of this same configuration agree with one another.

\section{Run-to-Run Variability and Reproducibility}\label{sec:reproducibility-audit}

This section reports the reproducibility audit design of the Brain-TokenGT arm
characterised in \autoref{sec:braintokengt}, then states the reproducibility property of
the specific implementation evaluated for this thesis's own architecture.

A single seed and a single run cannot distinguish a seed effect from ordinary run-to-run
scatter, so the stabilised Brain-TokenGT configuration was audited across nineteen repeat
runs without changing any setting, four to seven repeats at each of four seeds. The empirical
held-out test \ac{AUC} distributions, within-seed standard deviations, and head-to-head
comparisons are reported in \autoref{sec:rq3} (\autoref{tab:btgt-seed-variability}).

A run-to-run difference in non-deterministic models corresponds to a large number of pairwise
disagreements in the ranking the \ac{AUC} statistic summarises: at the test-set size used
for this comparator, a change of one hundredth of an \ac{AUC} point corresponds to
inverting the relative ranking of a handful of converter and stable-\ac{MCI} pairs, so a
difference on the order of 0.10 reflects dozens of such inversions between two runs that
share every setting, including the seed.

The mechanism consistent with this scatter is that determinism is enforced only at the
level of convolution algorithm selection, while the global deterministic-algorithm mode
that would additionally fix the order of floating-point reduction is never engaged.
Consequently, the scatter-and-gather operations inside the model's top-k pooling layer and
inside the recurrent core's node-embedding gather remain order-dependent on the graphics
processing unit, and their backward pass can accumulate floating-point error differently
between two runs that differ only in scheduling. This mechanism is consistent with the
observed scatter rather than demonstrated by an isolating experiment; no experiment in this
thesis holds the scheduling fixed while varying only the algorithm path. Engaging the
global deterministic-algorithm mode for this comparator was considered and set aside: doing
so raises an error inside operations that are part of the authors' released implementation
and would require rewriting them, which falls outside the scope of reproducing a published
architecture as released. The same pattern of held-out spread appears when the comparator
is evaluated on the pooled protocol described in \autoref{sec:pooled-protocol}, so the
scatter observed under the \ac{DELCODE}-only protocol is not specific to that cohort.

\subsection{The Trajectory Implementation Used in This Thesis Is Bit-Reproducible}\label{sec:bit-reproducible}

The temporal-first ladder introduced in \autoref{sec:tfgn} runs under the global
deterministic-algorithm mode by construction, the setting that \autoref{sec:reproducibility-audit}
identifies as absent from the Brain-TokenGT code path. Under that setting, repeating a run
of the \ac{GELSTM} none-arm implementation evaluated in this thesis at a fixed seed
reproduces its held-out test \ac{AUC} and cross-validation \ac{AUC} to the last recorded
digit. This was checked twice: once immediately after implementing the arm, and again after
a later change to how visit sequences are constructed for the pooled cohort, which
reproduced the earlier run's metrics exactly and confirmed that the change had not
introduced any nondeterminism of its own. This bit-reproducibility is a property of the
\ac{GELSTM} none-arm implementation evaluated here under this protocol; it is not stated of
any other arm or of an inherited pipeline, and it does not extend to the fine-tuned
spatial-first arm, whose seed-level reproducibility properties are reported alongside the
encoder ablation in \autoref{sec:rq2}.

\section{Comparability of the Head-to-Head}\label{sec:comparability}

Several elements of the comparison between Brain-TokenGT and the models introduced in this
thesis are held identical by construction. Both are built from the same
subject-level dataset object, so they see the same connectivity matrices for the same
subject set; both use the same split definitions and the same grouped five-fold
cross-validation; both select a decision threshold from out-of-fold predictions under the
same policy and apply it unchanged to the held-out test set; both are scored with the same
metric definitions; and both write their results into the same run-artefact schema. Four
further axes cannot be equalised in the same way, and each has a distinct consequence for
which comparisons the data support.

\textbf{The first axis is cohort windowing}. Brain-TokenGT is architecturally constrained to
sequences of two or three visits, since its published design targets cohorts with a fixed
number of acquisitions, whereas the trajectory models introduced in this thesis consume
sequences of one to six visits. Restricting to the two-to-three-visit window required by
Brain-TokenGT drops every subject with a single scan entirely, so a comparison that reads
one model at its full window and the other at the restricted window is reading the two
models on different subject counts as well as different architectures. This axis is
resolvable: \autoref{sec:rq3} reports a matched-window arm in which the trajectory models
are re-run restricted to the same two-to-three-visit window, holding the subject pool fixed
across both architectures. Brain-TokenGT itself is not re-run for this comparison, since it
already operates at exactly that window by construction, and restricting the trajectory
models to the same window was verified to remove visits without removing any subject from
the pool.

\textbf{The second axis is model capacity and tokenisation overhead}, which is not resolvable
without changing the published architecture. At two hundred regions and three visits,
Brain-TokenGT's token sequence is on the order of one thousand four hundred tokens per
subject, processed by a model with 603{,}849 trainable parameters, against a cross-validation
pool on the order of eighty training subjects per fold. The \ac{GELSTM} none arm processes
one pooled vector per visit with 31{,}169 trainable parameters, and its frozen-encoder
counterpart processes a sixty-four-dimensional latent vector per visit with 520{,}905
trainable parameters. This difference is a property of each published method rather than a
setting chosen for this comparison, so it cannot be equalised without evaluating a
different architecture than the one being compared. Where Brain-TokenGT's cross-validation
performance exceeds its held-out test performance by a wide margin, the observed
validation-to-test decrease is consistent with overfitting in this sample-size regime,
without asserting that overfitting is the only mechanism a token count this large relative
to the sample size could produce.

\textbf{The third axis is tuning budget.} The trajectory models introduced in this thesis were
developed and refined across the ablation ladder described in \autoref{sec:protocol},
whereas Brain-TokenGT was trained at its author-recommended configuration, adjusted only by
the learning-rate and weight-decay change described in \autoref{sec:braintokengt} that was
needed to obtain a configuration that trains to convergence at all. The published
architecture was originally tuned over two hundred trials by its authors; no comparably
sized search was run for either architecture under this protocol. This asymmetry is stated
without attributing any specific part of the observed performance difference to it, since
no experiment here isolates the tuning budget from the other three axes.

\textbf{The fourth axis is temporal representation.} The trajectory models introduced in this thesis
consume the interval between consecutive visits directly, whereas Brain-TokenGT treats
consecutive visits as uniform discrete steps with no notion of elapsed time. On
\ac{DELCODE}, this axis carries little weight in practice, since ninety percent of recorded
inter-visit intervals are twelve months apart; it is relevant chiefly on cohorts with
irregular intervals between visits.

Two consequences follow for which statistical comparisons are reported in
\autoref{sec:rq3}. First, because cohort windowing is the only one of the four axes that
can be equalised, the matched-window comparison is reported as a secondary comparison
rather than as the primary result: restricting both models to a window that discards the
visits the trajectory models exist to exploit cannot substitute for a comparison run at
each model's intended operating point. Second, because folds drawn from one seed share a
trained model and are not independent of one another, a comparison across the twenty folds
produced by five-fold cross-validation at four seeds is not treated as twenty independent
observations; the independent unit for any comparison against Brain-TokenGT is the seed, of
which four are available, and a rank-based test at that sample size is reported together
with its attainable significance floor rather than as a conventional significance test. No
conclusion is drawn from this comparison, or from any other result in this thesis, about
which of the two architectures would transfer to an external cohort: the run-to-run
variability documented in \autoref{sec:reproducibility-audit} is too large to support such
a conclusion, and the cross-cohort transfer results in \autoref{sec:rq4} are reported
without reference to it.
